\note{The authors also compared GA against three other trace-based methods: \textbf{Random Generation (RG)}, which samples traces uniformly at random as a baseline; \textbf{Bayesian Optimization (BO)} with a random forest surrogate model; and \textbf{Epsilon-Greedy Search (EPS)}, which mutates top-performing traces with probability $1-\epsilon$ and explores randomly with probability $\epsilon$.\par \par In Figure 3a, under a one-hour budget across 20 protocol pairs, GA consistently outperforms random generation across all three throughput weights $t_{\mathrm{coeff}} = 0, 0.5,$ and $1$ (improving scores by 23\% on average), supporting the benefit of using execution feedback under these tested settings. In Figure 3b, under a two-hour budget, GA, Bayesian optimization, and epsilon-greedy search achieve fairly similar performance, with GA having a slightly higher average, so the authors adopt GA as the default.
\par\smallskip Source: AdvNet paper §4.2.3, Fig.3}
